% appendix/questions/ordering.tex
% mainfile: ../../perfbook.tex
% SPDX-License-Identifier: CC-BY-SA-3.0

\section{需要多少排序？}
\label{sec:app:questions:How Much Ordering Is Needed?}
%
\epigraph{一寸黄金可以找回，一寸光阴永不再来。}
	 {中国谚语}

也许你精心构建了一个强排序的并发系统，
却发现它既不能良好地执行也不能良好地扩展。
或者你抛开谨慎，却发现你那卓越的快速和可扩展软件也不可靠。
是否存在一种既具有强健可靠性，
又具有强大性能和出色可扩展性的折中方案？

答案，一如既往，是"视情况而定"。

一种方法是先构建一个强排序系统，
然后检查其性能和可扩展性。
如果满足要求，系统就足够好了，不需要做更多的事情。
否则，进行仔细的分析
（参见\cref{sec:debugging:Performance Estimation}），
逐一解决每个瓶颈，直到系统的性能达标。

这种方法可以工作得很好，特别是与那种
对系统的随机组件进行优化以期望获得显著系统级收益的
常见做法相比。
然而，从强排序开始也可能相当浪费，
因为弱化系统瓶颈的排序可能需要大量重新设计和重写
系统其余部分以适应这种弱化。
更糟糕的是，消除一个瓶颈通常会暴露另一个瓶颈，
而后者又需要被弱化，进而可能导致系统其他部分的大规模重新设计和重写。
也许更糟糕的是另一种同样常见的做法，
即从一个快速但不可靠的系统开始，
然后无休止地与并发bug打地鼠，尽管在后一种情况下，
\cref{chp:Validation,chp:Formal Verification}
总是在那里为你提供帮助。

更好的做法是在设计时就有工具来确定系统的哪些部分
可以使用弱排序，同时确定哪些部分实际上从弱排序中受益。
以下各节将承担这些任务。

\subsection{定义数据在哪里？}
\label{sec:app:questions:Where is the Defining Data?}

实现这一目标的一种方法是牢记，
强排序所建立的一致性区域不能超出系统的边界。\footnote{
	这很可能是一个分布式系统。}
系统中那些角色是跟踪外部世界状态的部分
通常可以使用弱排序，因为光速延迟将迫使系统内的状态
落后于外部世界的状态。
花费大量开销来强制对本质上已经过时的数据保持一致的视图，
通常没有意义。
在这些情况下，\cref{chp:Deferred Processing}的方法
可能非常有帮助，\cref{chp:Data Structures}中描述的
一些数据结构也是如此。

尽管如此，采用某种对数据访问者来说有意义的语义是明智的，
例如，给定函数的返回值可能是：

\begin{enumerate}
\item	函数调用时的概念值与函数返回时的概念值之间的某个值。
	例如，参见\cref{sec:count:Statistical Counters}
	中讨论的统计计数器，记住这些计数器通常是单调的，
	至少在连续溢出之间是这样。
\item	函数调用到返回之间某个时刻的实际值。
	例如，参见\cref{lst:count:Just Count Atomically!}中
	所示的单变量原子计数器。
\item	如果该函数使用的值在函数调用和返回之间保持不变，
	则为预期值，否则为预期值的某种近似。
	精确指定近似的边界可能相当有挑战性。
	例如，考虑一个从RCU保护的链式数据结构的不同元素
	组合值的函数，如\cref{sec:datastruct:Read-Mostly Data Structures}
	中所述。
\end{enumerate}

较弱的排序通常意味着较弱的语义，
你应该能够向用户承诺这种弱化如何影响他们。
同时，除非调用者在函数调用和使用该函数计算的任何值期间
都持有锁，否则即使是完全有序的实现通常也不能
比上述选项给出的语义做得更好。

\QuickQuiz{
	但如果完全有序的实现不能提供比性能更好、
	可扩展性更强的弱排序实现更强的保证，
	那为什么还要费心使用完全排序？
}\QuickQuizAnswer{
	因为强排序实现有时能够在对给定数据结构的
	函数调用集合之间提供更大的一致性。
	例如，比较\cref{lst:count:Just Count Atomically!}的原子计数器
	和\cref{sec:count:Statistical Counters}的统计计数器。
	假设一个线程正在加3，另一个线程正在加5，
	同时另外两个线程正在并发读取计数器的值。
	使用原子计数器，不可能出现一个读者获得值3
	而另一个获得值5的情况。
	使用统计计数器，这种结果确实可能发生。
	事实上，在某些计算环境中，即使在像x86这样
	相对强排序的硬件上，这种结果也可能发生。

	因此，如果你的用户恰好需要这种公认不常见的一致性级别，
	你应该避免使用弱排序统计计数器。
}\QuickQuizEnd

有些人可能会争辩说，有用的计算只处理外部世界，
因此所有计算都可以使用弱排序。
这种论点是不正确的。
例如，你的银行账户余额是在银行的计算机内部定义的，
人们通常更喜欢涉及其账户余额的精确计算，
特别是那些可能怀疑任何此类近似都有利于银行的人。

简言之，虽然跟踪外部状态的数据可能是弱排序访问的
有吸引力的候选者，但请仔细考虑到底在跟踪什么以及什么在做跟踪。

\subsection{一致数据被一致使用了吗？}
\label{sec:app:questions:Consistent Data Used Consistently?}

弱化是安全的另一个暗示可能以这样的形式出现：
数据在持有锁时计算，但在释放锁后使用。
计算结果显然在锁释放后最多只是一个近似值，
这表明首先可以计算一个近似结果，
可能允许使用更弱的排序。
为此，\cref{chp:Counting}涵盖了许多近似计数方法。

然而，需要非常谨慎。
在锁释放后使用数据是弱排序优化可能有帮助的暗示？
还是锁释放过早的bug？

\subsection{问题可以分区吗？}
\label{sec:app:questions:Is the Problem Partitionable?}

假设系统持有数据的定义实例，
或者在锁释放后使用计算值被证明是一个bug。
那该怎么办？

一种方法是对系统进行分区，
如\cref{chp:Partitioning and Synchronization Design}中所讨论的。
分区可以提供出色的可扩展性，
在其更极端的形式中，可以实现与串行程序相媲美的每CPU性能，
如\cref{chp:Data Ownership}中所讨论的。
部分分区通常由锁来协调，
这是\cref{chp:Locking}的主题。

\subsection{以上皆不适用？}
\label{sec:app:questions:None of the Above?}

前面各节描述了获得性能和可扩展性的较简单方法，
有时使用更弱的排序，有时不使用。
但事实是，多核系统没有义务让生活变得简单。
也许\cref{sec:advsync:Advanced Synchronization,%
chp:Advanced Synchronization: Memory Ordering}
中涵盖的高级主题会有所帮助。

但请谨慎行事，因为优化非瓶颈代码来破坏代码库的稳定性太容易了。
再次强调，\cref{sec:debugging:Performance Estimation}可以提供帮助。
花时间回顾本书的其他部分也是值得的，
因为它包含了处理许多棘手情况的大量信息。
